% =====================================================================
%  Lecture 24 -- Control Plane                            (Persian edition)
% =====================================================================
\documentclass[11pt,a4paper]{article}
\providecommand\ACNroot{../..}
\input{\ACNroot/templates/lecture-fa}

\ACNset{number}{24}
\ACNset{title}{صفحهٔ کنترل نسل پنجم}
\ACNset{shorttitle}{صفحهٔ کنترل}
\ACNset{titleother}{The 5G Control Plane}
\ACNset{part}{بخش پنجم -- صفحه‌ها و چشم‌انداز (درس‌های ۲۳ تا ۲۶)}
\ACNset{time}{۶ تا ۸ ساعت (۳ ساعت مطالعه، ۳ تا ۵ ساعت تمرین)}
\ACNset{level}{پیشرفته}
\ACNset{prereq}{درس‌های ۴، ۱۲، ۱۸، ۲۰ و ۲۳}
\ACNset{gen}{5}

\begin{document}
\ACNcover

% ---------------------------------------------------------------------
\section{اهداف یادگیری}
پس از مطالعهٔ این درس باید بتوانید:
\begin{objectives}
  \item پشتهٔ پروتکل صفحهٔ کنترل نسل پنجم (\acro{RRC}، \acro{NAS}، \acro{NGAP}، \acro{SBI}) را \textbf{رسم کنید} و \textbf{توضیح دهید} هر پروتکل در کدام
        موجودیت‌ها خاتمه می‌یابد.
  \item حالت‌های \acro{RRC} یعنی \lr{IDLE}، \lr{INACTIVE} و \lr{CONNECTED} و رابطهٔ آن‌ها با حالت‌های \acro{CM} و \acro{RM} در هسته را \textbf{شرح دهید}.
  \item یک درخواست خدمت از \lr{CM-IDLE} را \textbf{دنبال کنید} و تأخیر صفحهٔ کنترل در گذارهای حالت را \textbf{برآورد کنید}.
  \item بار سیگنالینگ روی یک \acro{AMF}، معاوضهٔ زمان‌سنج بی‌کاری و اندازهٔ بهینهٔ ناحیهٔ ثبت را \textbf{محاسبه کنید}.
  \item \textbf{توضیح دهید} نسل پنجم چگونه صفحهٔ کنترل خود را در برابر اضافه‌بار و طوفان سیگنالینگ محافظت می‌کند.
\end{objectives}

% ---------------------------------------------------------------------
\section{دانش پیش‌نیاز}
به مدیریت مکان و فراخوانی (درس ۴)، \acro{MME} و کنترل اضافه‌بار آن (درس ۱۲)، معماری خدمت‌محور (درس ۱۸)، رویه‌های ثبت و نشست \acro{PDU} (درس
۲۰) و صفحهٔ کاربر درس ۲۳ نیاز دارید. این درس همان سامانه را از سوی سیگنالینگ می‌نگرد: چه کسی با چه کسی، در کدام حالت و با چه هزینه‌ای
گفت‌وگو می‌کند.

\paragraph{واژه‌های کلیدی} واژگان این درس همراه با معادل انگلیسی آن‌ها:
\begin{keyterms}
صفحهٔ کنترل & \lr{Control plane} & سیگنالینگی که اتصال را برقرار، نگهداری و رها می‌کند. \\
کنترل منابع رادیویی & \lr{RRC} & کنترل منابع رادیویی میان \acro{UE} و \acro{gNB}. \\
لایهٔ غیردسترسی & \lr{NAS (MM / SM)} & لایهٔ غیردسترسی میان \acro{UE} و \acro{AMF} (جابه‌جایی) یا \acro{SMF} (نشست‌ها). \\
پروتکل کاربرد \acro{NG} & \lr{NGAP} & پروتکل کاربرد روی \lr{N2} میان \acro{gNB} و \acro{AMF}، روی \acro{SCTP}. \\
حالت غیرفعال \acro{RRC} & \lr{RRC\_INACTIVE} & حالتی که زمینهٔ \acro{UE} حفظ می‌شود اما اتصال رادیویی معلق است. \\
حالت مدیریت اتصال / ثبت & \lr{CM / RM state} & دید هسته از \acro{UE}: متصل یا بیکار، ثبت‌شده یا ثبت‌نشده. \\
درخواست خدمت & \lr{Service request} & رویهٔ \acro{NAS} که \acro{UE} را از \lr{CM-IDLE} به \lr{CM-CONNECTED} می‌برد. \\
ناحیهٔ ثبت & \lr{Registration area} & فهرست نواحی ردیابی که \acro{UE} بیکار در آن‌ها لازم نیست مکان خود را گزارش کند. \\
ناحیهٔ اعلان مبتنی بر \acro{RAN} & \lr{RNA} & ناحیه‌ای که \acro{UE} غیرفعال بدون اطلاع به شبکه در آن جابه‌جا می‌شود. \\
طوفان سیگنالینگ & \lr{Signalling storm} & انفجاری از رویه‌های هم‌زمان که گره‌های صفحهٔ کنترل را دچار اضافه‌بار می‌کند. \\
\end{keyterms}

% ---------------------------------------------------------------------
\section{مقدمه و انگیزه}
صفحهٔ کنترل حتی یک بیت کاربر را حمل نمی‌کند، اما تعیین می‌کند آیا بیت‌های کاربر اصلاً می‌توانند جریان یابند یا نه. هر بار که یک تلفن هوشمند
برای دریافت یک اعلان بیدار می‌شود، یک تبادل پروتکلی کوچک در رادیو، \acro{gNB}، \acro{AMF}، \acro{SMF} و \acro{UPF} اجرا می‌کند. ضرب این در
میلیون‌ها دستگاه و ده‌ها گذار در ساعت، این سیگنالینگ بار پردازشی هسته و سهم چشمگیری از منابع رادیویی را تعیین می‌کند -- و کاربر تأخیر آن را
به‌صورت انتظار پیش از نخستین بسته حس می‌کند.

نسل پنجم صفحهٔ کنترل را از سه راه بازسازمان‌دهی می‌کند: یک حالت \acro{RRC} جدید، \lr{RRC\_INACTIVE}، که بیشتر سیگنالینگ هسته را برای انفجارهای
کوتاه حذف می‌کند؛ جداسازی روشن مدیریت جابه‌جایی (\acro{AMF}) و مدیریت نشست (\acro{SMF})؛ و واسط‌های خدمت‌محور درون هسته
\cite{ts23501,ts38300}. این درس پشتهٔ پروتکل و ماشین‌های حالت حاصل را توضیح می‌دهد و مدل‌های لازم برای ابعادگذاری آن‌ها را ارائه می‌کند.

% ---------------------------------------------------------------------
\section{مفاهیم بنیادی}

\subsection{پشتهٔ صفحهٔ کنترل}
سه خانوادهٔ پروتکل صفحهٔ کنترل را می‌سازند (\cref{fig:cp}):
\begin{itemize}
  \item \textbf{\acro{RRC}} میان \acro{UE} و \acro{gNB} \cite{ts38331}: اطلاعات سامانه، برقراری اتصال، فعال‌سازی امنیت رادیویی، پیکربندی حامل‌های
        رادیویی، اندازه‌گیری‌ها و فرمان‌های دست‌به‌دست. روی \acro{PDCP}، \acro{RLC} و \acro{MAC} در حامل‌های رادیویی سیگنالینگ اجرا می‌شود.
  \item \textbf{\acro{NAS}} میان \acro{UE} و هسته \cite{ts24501}: \lr{NAS-MM} (ثبت، درخواست خدمت، احراز هویت، امنیت \acro{NAS}) در \acro{AMF} خاتمه
        می‌یابد؛ \lr{NAS-SM} (برقراری و اصلاح نشست \acro{PDU}) توسط \acro{AMF} به \acro{SMF} بازپخش می‌شود. \acro{gNB} پیام‌های \acro{NAS} را بدون
        خواندن آن‌ها هدایت می‌کند -- این پیام‌ها سرتاسری میان \acro{UE} و \acro{AMF} حفاظت یکپارچگی و رمزنگاری می‌شوند.
  \item \textbf{\acro{NGAP}} میان \acro{gNB} و \acro{AMF} روی \lr{N2} \cite{ts38413}: مدیریت زمینهٔ \acro{UE}، برقراری منابع نشست \acro{PDU}، فراخوانی و
        سیگنالینگ دست‌به‌دست، روی \acro{SCTP}. درون هسته، \acro{AMF} و \acro{SMF} از واسط‌های خدمت‌محور (\lr{HTTP/2} با \acro{JSON}، درس ۱۸) استفاده
        می‌کنند.
\end{itemize}

\begin{acnfigure}
  \resizebox{\linewidth}{!}{\acnfig{lecture-24/cp-stack}}
  \caption{پشتهٔ پروتکل صفحهٔ کنترل نسل پنجم: \acro{RRC} روی \lr{Uu}، \acro{NGAP} روی \lr{N2}، \acro{NAS} سرتاسری میان \acro{UE} و \acro{AMF}/\acro{SMF}، و
  \acro{SBI} در هسته.}
  \label{fig:cp}
\end{acnfigure}

\subsection{حالت‌های یک \lr{UE}}
شبکه برای هر \acro{UE} سه ماشین حالت جفت‌شده نگه می‌دارد \cite{ts23501,ts38331}:
\begin{itemize}
  \item \textbf{\acro{RM}} (مدیریت ثبت) در \acro{AMF}: \lr{RM-DEREGISTERED} یا \lr{RM-REGISTERED}.
  \item \textbf{\acro{CM}} (مدیریت اتصال) در \acro{AMF}: \lr{CM-IDLE} (بدون اتصال \lr{N2} برای \acro{UE}؛ مکان در حد ناحیهٔ ثبت معلوم است) یا
        \lr{CM-CONNECTED} (یک اتصال \lr{N2} وجود دارد).
  \item \textbf{\acro{RRC}} در \acro{gNB} و \acro{UE}: \lr{RRC\_IDLE}، \lr{RRC\_INACTIVE} یا \lr{RRC\_CONNECTED} (\cref{fig:rrc}).
\end{itemize}
در \lr{RRC\_INACTIVE}، \acro{UE} تقریباً مانند یک \acro{UE} بیکار رفتار می‌کند -- روی سلول‌ها مستقر می‌شود، فراخوانی را پایش می‌کند و درون یک ناحیهٔ
اعلان مبتنی بر \acro{RAN} (\acro{RNA}) آزادانه جابه‌جا می‌شود -- اما \acro{gNB}ای که آن را معلق کرده زمینهٔ \acro{UE} را نگه می‌دارد و اتصال‌های
\lr{N2}/\lr{N3} به هسته برقرار می‌مانند، پس هسته \acro{UE} را \lr{CM-CONNECTED} می‌بیند. بازگشت به \lr{RRC\_CONNECTED} تنها به یک تبادل
\lr{RRC Resume} بدون سیگنالینگ هسته نیاز دارد.

\begin{acnfigure}
  \resizebox{0.92\linewidth}{!}{\acnfig{lecture-24/rrc-states}}
  \caption{حالت‌ها و گذارهای \acro{RRC} در \acro{NR}؛ \lr{RRC\_INACTIVE} در هسته به \lr{CM-CONNECTED} تعلق دارد.}
  \label{fig:rrc}
\end{acnfigure}

\subsection{دسترس‌پذیری}
اینکه شبکه \acro{UE} را کجا می‌یابد به حالت آن بستگی دارد: در \lr{CM-IDLE}، \acro{AMF} ناحیهٔ ثبت (فهرستی از نواحی ردیابی) را می‌داند و باید در آن
فراخوانی کند؛ در \lr{RRC\_INACTIVE}، \acro{gNB} لنگر \acro{RNA} را می‌داند و خودش فراخوانی \acro{RAN} را انجام می‌دهد؛ در \lr{RRC\_CONNECTED} سلول
خدمت‌دهنده معلوم است. \acro{UE} خروج از ناحیهٔ خود را با یک به‌روزرسانی ثبت جابه‌جایی (\lr{CM-IDLE}) یا یک به‌روزرسانی \acro{RNA}
(\lr{RRC\_INACTIVE}) گزارش و آن را به‌طور دوره‌ای تازه می‌کند.

% ---------------------------------------------------------------------
\section{توضیحات فنی تفصیلی}

\subsection{درخواست خدمت}
وقتی یک \acro{UE} بیکار دادهٔ فراسو دارد -- یا برای دادهٔ فروسو فراخوانده می‌شود -- درخواست خدمت را اجرا می‌کند (\cref{fig:sr}) \cite{ts23502}. پس
از دسترسی تصادفی، \acro{UE} یک اتصال \acro{RRC} برقرار می‌کند و \lr{NAS Service Request} را همراه \lr{RRC Setup Complete} می‌فرستد؛ \acro{gNB} آن را
در یک \lr{NGAP Initial UE Message} به \acro{AMF}ای که با \lr{5G-S-TMSI} انتخاب شده هدایت می‌کند. \acro{AMF} یکپارچگی \acro{NAS} را بررسی می‌کند،
از \acro{SMF} هر نشست \acro{PDU} که باید فعال شود می‌خواهد صفحهٔ کاربرش را دوباره برقرار کند (\acro{SMF} \acro{UPF} را روی \lr{N4} به‌روز می‌کند)،
سپس یک \lr{Initial Context Setup Request} حاوی کلید امنیتی و منابع نشست می‌فرستد. \acro{gNB} امنیت \acro{AS} را فعال، حامل‌های رادیویی را دوباره
پیکربندی و نقاط پایانی تونل فروسوی خود را برمی‌گرداند که \acro{SMF} آن‌ها را در \acro{UPF} نصب می‌کند.

\begin{acnfigure}
  \resizebox{\linewidth}{!}{\acnfig{lecture-24/service-request}}
  \caption{درخواست خدمت آغازشده از سوی \acro{UE} از \lr{CM-IDLE} (ساده‌شده، بر اساس \lr{TS 23.502}).}
  \label{fig:sr}
\end{acnfigure}

\subsection{تعلیق و از سرگیری}
با \lr{RRC\_INACTIVE}، \acro{gNB} یک اتصال را با یک \lr{RRC Release} همراه با نشانگر تعلیق پایان می‌دهد: \acro{UE} و \acro{gNB} زمینهٔ \acro{AS} و پیکربندی
امنیتی را نگه می‌دارند و \acro{UE} یک شناسهٔ از سرگیری و یک \acro{RNA} دریافت می‌کند. وقتی \acro{UE} دوباره داده دارد، یک \lr{RRC Resume Request}
می‌فرستد؛ اگر در سلولی از \acro{gNB} دیگری باشد، آن \acro{gNB} زمینه را روی \lr{Xn} از \acro{gNB} لنگر بازیابی می‌کند و از \acro{AMF} می‌خواهد مسیر
\lr{N2} را تغییر دهد. در حالت رایج از سرگیری در همان ناحیه، هیچ سیگنالینگ \acro{NAS} یا هسته‌ای لازم نیست. مقادیر کوچک داده را حتی می‌توان در
خود رویهٔ از سرگیری فرستاد (ارسال دادهٔ کوچک، انتشار ۱۷).

\subsection{جداسازی مدیریت جابه‌جایی و نشست}
در \acro{EPC}، \acro{MME} هم سیگنالینگ جابه‌جایی و هم نشست را رسیدگی می‌کرد (درس ۱۲). در نسل پنجم \acro{AMF} تنها جابه‌جایی، دسترسی و امنیت \acro{NAS}
را رسیدگی می‌کند و \acro{SMF} نشست‌ها را. یک \acro{UE} یک \acro{AMF} دارد اما ممکن است چند \acro{SMF} داشته باشد (یکی برای هر نشست \acro{PDU} یا
برش). این جداسازی اجازه می‌دهد دسترسی غیر\lr{3GPP} (\lr{Wi-Fi} از طریق \acro{N3IWF}) همان \acro{AMF} را به اشتراک بگذارد، \acro{AMF} را از انواع نشست
مستقل نگه می‌دارد و به برش‌ها اجازه می‌دهد \acro{SMF}های خود را داشته باشند (درس ۲۱).

\subsection{کنترل اضافه‌بار و ازدحام}
گره‌های صفحهٔ کنترل در چند سطح از خود محافظت می‌کنند \cite{ts23501,ts24501}: کنترل دسترسی یکپارچه در \acro{UE} دسته‌های دسترسی‌ای را که \acro{gNB} پخش
می‌کند پیش از ارسال هر پیامی مسدود می‌کند؛ \acro{AMF} می‌تواند از \acro{gNB}ها بخواهد اتصال‌های \acro{RRC} جدید را رد یا کم کنند (\lr{NGAP overload
start})؛ کنترل ازدحام \acro{NAS} در پیام‌های رد یک زمان‌سنج عقب‌نشینی (مثلاً \lr{T3346}) برمی‌گرداند تا \acro{UE}ها بعداً دوباره تلاش کنند؛ و \acro{NRF} و
\acro{SCP} بار را میان چند نمونهٔ \acro{AMF} و \acro{SMF} پخش می‌کنند (درس ۱۸).

% ---------------------------------------------------------------------
\section{مدل‌های ریاضی}

\subsection{بار سیگنالینگ}\label{sec:load}
\paragraph{مدل} اگر $N$ دستگاه \acro{UE} هر کدام رویهٔ $k$ را با نرخ $\lambda_k$ در ساعت انجام دهند و هر رویه $m_k$ پیام در \acro{AMF} هزینه داشته
باشد، نرخ پیام بر حسب پیام در ثانیه برابر است با
\begin{equation}\label{eq:load}
  \Lambda = \frac{N}{3600} \sum_{k} \lambda_k\, m_k .
\end{equation}
\paragraph{مثال} $N = 2\times10^{6}$ با نرخ‌های ساعت اوج برای هر \acro{UE}: ۲۰ درخواست خدمت (هر کدام ۴ پیام)، ۲ ثبت جابه‌جایی (۶)، ۳ دست‌به‌دست (۶)،
۰٫۱ ثبت اولیه (۲۰) و ۰٫۵ برقراری نشست \acro{PDU} (۸): $\sum_k \lambda_k m_k = 116$ پیام برای هر \acro{UE} در ساعت و $\Lambda = 64\,444$ پیام بر ثانیه. اگر
\lr{RRC\_INACTIVE} ۸۰٪ گذارهای بیکار به متصل را به از سرگیری بدون سیگنالینگ هسته تبدیل کند، مجموع به ۵۲ و $\Lambda = 28\,889$ پیام بر ثانیه کاهش
می‌یابد -- کاهشی ۵۵ درصدی.

\subsection{زمان‌سنج بی‌کاری}\label{sec:timer}
\paragraph{مدل} یک \acro{UE} انفجارهای داده را به‌صورت یک فرایند پواسون با نرخ $\lambda$ تولید می‌کند. پس از هر انفجار، \acro{gNB} اتصال را به مدت
بی‌کاری $T$ نگه می‌دارد و سپس آن را رها (یا معلق) می‌کند. گذار به بیکار پس از یک انفجار تنها وقتی رخ می‌دهد که انفجار بعدی دیرتر از $T$ بیاید،
پس نرخ گذار و میانگین «دنبالهٔ» متصل برای هر انفجار عبارت‌اند از
\begin{equation}\label{eq:timer}
  r(T) = \lambda\, e^{-\lambda T}, \qquad \bar{t}_{\mathrm{tail}}(T) = \frac{1 - e^{-\lambda T}}{\lambda} .
\end{equation}
کسری از زمان که در دنباله سپری می‌شود $\lambda\,\bar{t}_{\mathrm{tail}} = 1 - e^{-\lambda T}$ است.
\paragraph{مثال} یک انفجار در دقیقه ($\lambda = 1/60$\,s$^{-1}$): $T = 10$ ثانیه ۵۰٫۸ گذار در ساعت و ۱۵٫۴٪ زمان در دنباله می‌دهد؛ $T = 30$ ثانیه ۳۶٫۴
گذار در ساعت و ۳۹٫۳٪ (\cref{fig:timer}). زمان‌سنج بلندتر سیگنالینگ را صرفه‌جویی می‌کند اما انرژی \acro{UE} و منابع رادیویی را هزینه می‌کند؛
\lr{RRC\_INACTIVE} زمان‌سنج‌های کوتاه را مقرون‌به‌صرفه می‌کند چون هر گذار ارزان است.

\begin{acnfigure}
  \acnfig{lecture-24/inactivity-timer}
  \caption{معاوضهٔ زمان‌سنج بی‌کاری برای یک انفجار در دقیقه (\cref{eq:timer}).}
  \label{fig:timer}
\end{acnfigure}

\subsection{اندازهٔ ناحیهٔ ثبت}\label{sec:ra}
\paragraph{مدل} با مدل جابه‌جایی جریان سیال (درس ۴)، یک \acro{UE} که با سرعت میانگین $v$ در ناحیه‌ای با محیط $P$ و مساحت $A$ حرکت می‌کند با نرخ
$vP/(\pi A)$ از مرز آن عبور می‌کند. برای یک ناحیهٔ مربعی با ضلع $a$ این مقدار $R_u = 4v/(\pi a)$ است. اگر هر به‌روزرسانی $w_u$ هزینه داشته باشد و هر
رویداد فراخوانی (با نرخ $\lambda_p$) در هر یک از $n = a^2/A_c$ سلول با مساحت $A_c$ هزینهٔ $w_p$ داشته باشد، هزینه برای هر \acro{UE} در ساعت
\begin{equation}\label{eq:racost}
  C(a) = w_u \frac{4v}{\pi a} + w_p \lambda_p \frac{a^2}{A_c}
\end{equation}
است که با
\begin{equation}\label{eq:raopt}
  a^{*} = \left( \frac{2 w_u v A_c}{\pi w_p \lambda_p} \right)^{1/3}
\end{equation}
کمینه می‌شود. در نقطهٔ بهینه هزینهٔ به‌روزرسانی دقیقاً دو برابر هزینهٔ فراخوانی است.
\paragraph{مثال} $v = 20$ کیلومتر بر ساعت، $A_c = 1$ کیلومتر مربع، $\lambda_p = 5$ در ساعت، $w_u = 10$، $w_p = 1$: $a^{*} = 2.94$ کیلومتر (حدود ۸٫۷ سلول در هر
ناحیه)، با هزینهٔ به‌روزرسانی ۸۶٫۶ و هزینهٔ فراخوانی ۴۳٫۳ برای هر \acro{UE} در ساعت (\cref{fig:ra}).

\begin{acnfigure}
  \acnfig{lecture-24/ta-cost}
  \caption{معاوضهٔ اندازهٔ ناحیهٔ ثبت (\cref{eq:racost,eq:raopt}).}
  \label{fig:ra}
\end{acnfigure}

\subsection{تأخیر صفحهٔ کنترل}\label{sec:cplat}
\paragraph{مدل} زمان از یک حالت کم‌مصرف تا آغاز انتقال داده مجموع گام‌های گذار است:
\begin{equation}\label{eq:cplat}
  T_{\mathrm{CP}} = T_{\mathrm{RA}} + T_{\mathrm{RRC}} + T_{\mathrm{core}} + T_{\mathrm{ctx}} .
\end{equation}
\lr{TR 38.913} هدف ۱۰ میلی‌ثانیه را تعیین می‌کند \cite{tr38913}؛ \lr{ITU-R M.2410} ۲۰ میلی‌ثانیه را الزامی و ۱۰ میلی‌ثانیه را توصیه می‌کند \cite{itu2017m2410}.
\paragraph{مثال} با مقادیر نمونه برای فاصلهٔ زیرحامل ۳۰ کیلوهرتز -- دسترسی تصادفی ۲ میلی‌ثانیه، تبادل \acro{RRC} ۴ میلی‌ثانیه، سیگنالینگ هسته
(\lr{N2}، \acro{AMF}، \acro{SMF} و \lr{N4}) ۶ میلی‌ثانیه، برقراری زمینه با امنیت \acro{AS} و پیکربندی دوباره ۴ میلی‌ثانیه -- درخواست خدمت از
\lr{RRC\_IDLE} ۱۶ میلی‌ثانیه طول می‌کشد. از سرگیری از \lr{RRC\_INACTIVE} تنها به دسترسی تصادفی و تبادل \acro{RRC} نیاز دارد: ۶ میلی‌ثانیه. تنها
حالت غیرفعال هدف ۱۰ میلی‌ثانیه را برآورده می‌کند.

% ---------------------------------------------------------------------
\section{معماری شبکه و نمودارها}
صفحهٔ کنترل \acro{UE}، \acro{gNB} (پس از تقسیم، \lr{CU-CP}؛ درس ۲۳)، مجموعهٔ \acro{AMF}ای که به یک منطقه خدمت می‌دهد، و \acro{SMF}ها، \acro{PCF}،
\acro{UDM} و \acro{AUSF} را که از طریق \acro{SBA} در دسترس‌اند در بر می‌گیرد. \acro{AMF}ها به‌صورت مجموعه‌هایی از نمونه‌های جایگزین‌پذیر مستقر
می‌شوند که زمینه‌های \acro{UE} را در یک تابع ذخیره‌سازی به اشتراک می‌گذارند تا یک \acro{gNB} بتواند به هر نمونه از مجموعه برسد (یک \acro{GUAMI}
مجموعه و اشاره‌گر را مشخص می‌کند). \cref{fig:cp,fig:rrc,fig:sr} دید پروتکلی، دید حالت و دید رویه‌ای این معماری را نشان می‌دهند.

% ---------------------------------------------------------------------
\section{رویه‌های گام‌به‌گام}

\begin{procedurebox}{دادهٔ پایان‌یافته در موبایل برای یک \lr{UE} بیکار}\label{proc:mt}
\begin{steps}
  \item یک بستهٔ فروسو برای نشستی بدون تونل \lr{N3} به \acro{UPF} می‌رسد؛ \acro{UPF} آن را بافر می‌کند (یا به \acro{SMF} می‌فرستد) و به \acro{SMF} اطلاع می‌دهد.
  \item \acro{SMF} از \acro{AMF} می‌خواهد \acro{UE} را دسترس‌پذیر کند (\lr{Namf\_Communication\_N1N2MessageTransfer}).
  \item \acro{AMF} \acro{UE} را در همهٔ سلول‌های ناحیهٔ ثبت آن از طریق \acro{gNB}ها فرامی‌خواند (\lr{NGAP Paging}).
  \item \acro{UE} با یک درخواست خدمت پاسخ می‌دهد (\cref{fig:sr})؛ صفحهٔ کاربر دوباره برقرار و دادهٔ بافرشده تحویل می‌شود.
\end{steps}
\end{procedurebox}

\begin{procedurebox}{از سرگیری از \lr{RRC\_INACTIVE}}\label{proc:resume}
\begin{steps}
  \item \acro{UE} (با دادهٔ فراسو یا پس از فراخوانی \acro{RAN}) دسترسی تصادفی انجام می‌دهد و یک \lr{RRC Resume Request} با شناسهٔ از سرگیری خود می‌فرستد.
  \item اگر زمینه نزد \acro{gNB} دیگری باشد، \acro{gNB} جدید آن را روی \lr{Xn} بازیابی می‌کند و سپس از \acro{AMF} تغییر مسیر \lr{N2} را می‌خواهد.
  \item \acro{gNB} با \lr{RRC Resume} پاسخ می‌دهد؛ امنیت از زمینهٔ ذخیره‌شده با کلیدهای تازه بازگردانده می‌شود.
  \item \acro{UE} با \lr{RRC Resume Complete} تأیید می‌کند و داده روی تونل‌های \lr{N3} موجود جریان می‌یابد.
\end{steps}
\end{procedurebox}

% ---------------------------------------------------------------------
\section{مثال‌های حل‌شده}

\begin{examplebox}{ابعادگذاری یک مجموعهٔ \lr{AMF}}\label{ex:load}
\textbf{مسئله.} برای بار \cref{sec:load}، اگر هر نمونه $\mu = 20\,000$ پیام بر ثانیه پردازش کند و بهره‌وری آن نباید از ۷۰٪ بیشتر شود، چند نمونهٔ \acro{AMF} لازم
است؟ با و بدون \lr{RRC\_INACTIVE} مقایسه کنید.

\textbf{حل.} ظرفیت قابل استفاده برای هر نمونه $0.7\mu = 14\,000$ پیام بر ثانیه. بدون \lr{INACTIVE}: $64\,444/14\,000 = 4.6$، پس ۵ نمونه. با \lr{INACTIVE}:
$28\,889/14\,000 = 2.06$، پس ۳ نمونه. کارایی سیگنالینگ در \acro{RAN} مستقیماً به نمونه‌های کمتر در هسته تبدیل می‌شود.
\end{examplebox}

\begin{examplebox}{انتخاب زمان‌سنج بی‌کاری}\label{ex:timer}
\textbf{مسئله.} با یک انفجار در دقیقه، بلندترین زمان‌سنجی که \acro{UE} را حداکثر ۲۰٪ زمان در دنباله نگه دارد چیست و چند گذار در ساعت حاصل می‌شود؟

\textbf{حل.} $1 - e^{-T/60} \le 0.2 \Rightarrow T \le -60\ln 0.8 = 13.4$ ثانیه. گذارها: $60\,e^{-13.4/60} = 60\times0.8 = 48$ در ساعت.
\end{examplebox}

\begin{examplebox}{ناحیهٔ ثبت برای کاربران سریع}\label{ex:ra}
\textbf{مسئله.} مثال \cref{sec:ra} را برای کاربران خودرویی با $v = 60$ کیلومتر بر ساعت تکرار کنید.

\textbf{حل.} $a^{*3} = 2\times10\times60\times1/(\pi\times5) = 76.4$، $a^{*} = 4.24$ کیلومتر (حدود ۱۸ سلول). هزینهٔ به‌روزرسانی $10\times4\times60/(\pi\times4.24) = 180$،
هزینهٔ فراخوانی $5\times4.24^2 = 90$، مجموع ۲۷۰ برای هر \acro{UE} در ساعت. کاربران سریع‌تر به نواحی بزرگ‌تر نیاز دارند؛ به همین دلیل اپراتورها به
فهرست‌های \acro{TA} در امتداد بزرگراه‌ها و راه‌آهن سلول‌های بیشتری می‌دهند.
\end{examplebox}

% ---------------------------------------------------------------------
\section{مقایسهٔ فناوری‌ها}
\begin{table}[htbp]
\centering\small
\caption{صفحهٔ کنترل در \lr{LTE/EPC} و در \lr{NR/5GC}.}\label{tab:cmp}
\rowcolors{2}{acnmist}{white}
\begin{tabularx}{\linewidth}{@{}l Y Y@{}}
\toprule
\rowcolor{white} & \lr{LTE / EPC} & \lr{NR / 5GC} \\
\midrule
حالت‌های \acro{RRC}   & \lr{IDLE}، \lr{CONNECTED} (تعلیق/از سرگیری فقط برای اینترنت اشیا) & \lr{IDLE}، \lr{INACTIVE}، \lr{CONNECTED} \\
جابه‌جایی و نشست‌ها   & \acro{MME} (هر دو)، کنترل \lr{S-GW/P-GW} & \acro{AMF} (جابه‌جایی)، \acro{SMF} (نشست‌ها) \\
واسط‌های هسته        & \lr{Diameter}، \lr{GTPv2-C} (نقطه به نقطه) & \acro{SBI} روی \lr{HTTP/2}، کشف با \acro{NRF} \\
پروتکل \acro{RAN}--هسته & \lr{S1AP} روی \acro{SCTP} & \acro{NGAP} روی \acro{SCTP} \\
نواحی مکان           & فهرست \acro{TA}، فراخوانی توسط \acro{MME} & فهرست \acro{TA} (\lr{CM-IDLE}) و \acro{RNA} با فراخوانی \acro{RAN} (\lr{INACTIVE}) \\
کنترل دسترسی         & \lr{ACB}، \lr{EAB}، زمان‌سنج‌های عقب‌نشینی & کنترل دسترسی یکپارچه، زمان‌سنج‌های عقب‌نشینی \\
دسترسی غیر\lr{3GPP}   & مسیر جداگانهٔ \lr{ePDG} & همان \acro{AMF} از طریق \acro{N3IWF} \\
\bottomrule
\end{tabularx}
\end{table}

\paragraph{سیگنالینگ بهای کارایی است} هر دورهٔ بی‌کاری انرژی و منابع رادیویی را صرفه‌جویی می‌کند اما باید بهای آن را با یک گذار پرداخت. صفحهٔ کنترل
نسل پنجم این بها را کاهش می‌دهد (\lr{INACTIVE}، پیام‌های کمتر در هسته، پردازش سریع‌تر)، و این به اپراتورها اجازه می‌دهد اتصال‌ها را زودتر رها کنند
و به ازای هر واحد ظرفیت هسته به دستگاه‌های بسیار بیشتری خدمت دهند.

% ---------------------------------------------------------------------
\section{نکات کلیدی}
\begin{takeawaybox}
\begin{itemize}
  \item \acro{RRC} (\acro{UE}--\acro{gNB})، \lr{NAS-MM} (\acro{UE}--\acro{AMF})، \lr{NAS-SM} (\acro{UE}--\acro{SMF} از طریق \acro{AMF})، \acro{NGAP} (\acro{gNB}--\acro{AMF}
        روی \acro{SCTP}) و \acro{SBI} در هسته.
  \item \lr{RRC\_INACTIVE} زمینه را در \acro{RAN} و هسته را در \lr{CM-CONNECTED} نگه می‌دارد: از سرگیری به سیگنالینگ هسته نیاز ندارد.
  \item بار $\Lambda = (N/3600)\sum_k \lambda_k m_k$؛ گذارهای بیکار به متصل بر آن غالب‌اند.
  \item زمان‌سنج بی‌کاری: گذارها $\lambda e^{-\lambda T}$ در برابر کسر دنبالهٔ $1 - e^{-\lambda T}$.
  \item ناحیهٔ ثبت: $a^{*} = (2w_u v A_c/(\pi w_p\lambda_p))^{1/3}$؛ در نقطهٔ بهینه هزینهٔ به‌روزرسانی دو برابر فراخوانی است.
\end{itemize}
\end{takeawaybox}

% ---------------------------------------------------------------------
\section{تمرین‌ها}

\begin{exercisebox}{1}{\lr{UE} کجاست؟}\label{exr:states}
برای هر حالت \acro{RRC}، حالت \acro{CM} از دید \acro{AMF} و دقتی را که شبکه با آن مکان \acro{UE} را می‌داند بیان کنید.
\end{exercisebox}

\begin{exercisebox}{1}{\lr{NAS} شفاف}\label{exr:nas}
چرا \acro{NAS} برای \acro{gNB} شفاف است؟ یک دلیل امنیتی و یک دلیل معماری بیاورید.
\end{exercisebox}

\begin{exercisebox}{2}{بار \lr{AMF}}\label{exr:load}
یک مجموعهٔ \acro{AMF} به $5\times10^{5}$ دستگاه \acro{UE} با نرخ‌های ساعت اوج ۳۰ درخواست خدمت (هر کدام ۴ پیام)، ۱ ثبت جابه‌جایی (۶) و ۴ دست‌به‌دست (۶) برای هر
\acro{UE} خدمت می‌دهد. نرخ پیام را محاسبه کنید.
\end{exercisebox}

\begin{exercisebox}{2}{کاربردهای پرحرف}\label{exr:timer}
یک کاربرد پیام‌رسان به‌طور میانگین هر ۲۰ ثانیه یک انفجار تولید می‌کند. گذارها در ساعت و کسر دنباله را برای $T = 5$ و $T = 10$ ثانیه محاسبه کنید.
\end{exercisebox}

\begin{exercisebox}{3}{ناحیهٔ بهینه}\label{exr:ra}
\cref{eq:raopt} را از \cref{eq:racost} استخراج کنید، نشان دهید هزینهٔ به‌روزرسانی در نقطهٔ بهینه دو برابر هزینهٔ فراخوانی است، و $a^{*}$ را برای عابران پیاده
($v = 5$ کیلومتر بر ساعت) با سایر مقادیر \cref{sec:ra} محاسبه کنید.
\end{exercisebox}

\begin{exercisebox}{3}{طوفان سیگنالینگ}\label{exr:storm}
پس از یک قطعی هسته، $10^{6}$ دستگاه \acro{UE} باید دوباره ثبت شوند و هر ثبت ۲۰ پیام در یک مجموعهٔ \acro{AMF} با ۵ نمونه و $\mu = 20\,000$ پیام بر ثانیه هزینه
دارد. اگر \acro{UE}ها تلاش‌های خود را به‌طور یکنواخت در پنجرهٔ $W$ پخش کنند، کوچک‌ترین $W$ برای بهره‌وری ۱۰۰٪ و ۷۰٪ چیست؟ کدام سازوکارها چنین پخشی را
تحمیل می‌کنند؟
\end{exercisebox}

\subsection*{پاسخ‌ها}
\addcontentsline{toc}{subsection}{پاسخ‌ها}

\begin{solutionbox}{\ref{exr:states}}
\lr{RRC\_IDLE}: \lr{CM-IDLE}، ناحیهٔ ثبت (فهرست \acro{TA}ها). \lr{RRC\_INACTIVE}: \lr{CM-CONNECTED}، \acro{RNA} که \acro{gNB} لنگر می‌داند (\acro{AMF} \acro{gNB} را
می‌داند). \lr{RRC\_CONNECTED}: \lr{CM-CONNECTED}، سلول خدمت‌دهنده.
\end{solutionbox}

\begin{solutionbox}{\ref{exr:nas}}
امنیت: پیام‌های \acro{NAS} سرتاسری با کلیدهایی حفاظت می‌شوند که تنها \acro{UE} و \acro{AMF} در اختیار دارند، پس یک گره \acro{RAN} نفوذشده یا نامطمئن
(از جمله دسترسی غیر\lr{3GPP}) نمی‌تواند آن‌ها را بخواند یا تغییر دهد. معماری: \acro{RAN} از رویه‌های هسته مستقل است؛ قابلیت‌های جدید \acro{NAS} به تغییر
\acro{RAN} نیاز ندارند و همان \acro{NAS} روی دسترسی \lr{3GPP} و غیر\lr{3GPP} اجرا می‌شود.
\end{solutionbox}

\begin{solutionbox}{\ref{exr:load}}
$\sum_k \lambda_k m_k = 120 + 6 + 24 = 150$ برای هر \acro{UE} در ساعت؛ $\Lambda = 5\times10^5\times150/3600 = 20\,833$ پیام بر ثانیه.
\end{solutionbox}

\begin{solutionbox}{\ref{exr:timer}}
$\lambda = 1/20$\,s$^{-1}$، ۱۸۰ انفجار در ساعت. $T = 5$ ثانیه: $180\,e^{-0.25} = 140.2$ گذار در ساعت، دنبالهٔ $1 - e^{-0.25} = 22.1\,\%$. $T = 10$ ثانیه:
$180\,e^{-0.5} = 109.2$، دنبالهٔ ۳۹٫۳٪. انفجارهای کوچک و پرتکرار همان حالتی است که \lr{RRC\_INACTIVE} و ارسال دادهٔ کوچک برای آن طراحی شده‌اند.
\end{solutionbox}

\begin{solutionbox}{\ref{exr:ra}}
$C'(a) = -4w_uv/(\pi a^2) + 2w_p\lambda_p a/A_c = 0 \Rightarrow a^3 = 2w_uvA_c/(\pi w_p\lambda_p)$. در نقطهٔ بهینه $w_p\lambda_p a^2/A_c = 2w_uv/(\pi a)$، یعنی هزینهٔ
فراخوانی $= 2w_uv/(\pi a)$، در حالی که هزینهٔ به‌روزرسانی $= 4w_uv/(\pi a)$: دو برابر. عابران پیاده: $a^{*3} = 2\times10\times5/(5\pi) = 6.37$، $a^{*} = 1.85$ کیلومتر.
\end{solutionbox}

\begin{solutionbox}{\ref{exr:storm}}
مجموع $2\times10^{7}$ پیام؛ ظرفیت $10^{5}$ پیام بر ثانیه. ۱۰۰٪: $W \ge 200$ ثانیه؛ ۷۰٪: $W \ge 2\times10^{7}/(7\times10^{4}) = 286$ ثانیه. سازوکارها: زمان‌سنج‌های
عقب‌نشینی تصادفی در پیام‌های رد \acro{NAS} (مثلاً \lr{T3346})، مسدودسازی کنترل دسترسی یکپارچه که \acro{gNB}ها پخش می‌کنند، نشانگرهای اضافه‌بار \acro{AMF}
به \acro{RAN}، و اولویت‌دهی به دسترسی اضطراری و پراولویت.
\end{solutionbox}

% ---------------------------------------------------------------------
\section{مطالعهٔ بیشتر}
پروتکل و حالت‌های \acro{RRC} در \lr{TS 38.331} \cite{ts38331} مشخص و در \lr{TS 38.300} \cite{ts38300} خلاصه شده‌اند؛ \acro{NAS} در \lr{TS 24.501} \cite{ts24501}؛
\acro{NGAP} در \lr{TS 38.413} \cite{ts38413}؛ حالت‌های \acro{CM} و \acro{RM} در \lr{TS 23.501} \cite{ts23501} و رویه‌ها در \lr{TS 23.502} \cite{ts23502}. اهداف تأخیر از
\lr{TR 38.913} \cite{tr38913} و \lr{ITU-R M.2410} \cite{itu2017m2410} می‌آیند. رومر و همکاران \cite{rommer2020core} سیگنالینگ هستهٔ نسل پنجم را توضیح می‌دهند؛
دالمن و همکاران \cite{dahlman2020nr} حالت‌های \acro{RRC} و ارسال دادهٔ کوچک را.

\begin{latin}
\printbibliography[heading=subbibliography,title={References}]
\end{latin}
\end{document}
